\documentclass[12pt]{article}
\usepackage[utf8]{inputenc}
\usepackage[T1]{fontenc}
\usepackage[margin=1in]{geometry}
\usepackage{amsmath,amssymb}
\usepackage{array}
\usepackage{enumitem}
\usepackage{fancyhdr}
\usepackage{graphicx}
\usepackage{booktabs}
\usepackage{lastpage}

\pagestyle{fancy}
\fancyhf{}
\fancyhead[L]{\small VIETNAM NATIONAL UNIVERSITY -- HCMC \\ \small UNIVERSITY OF SCIENCE}
\fancyhead[R]{\small \textbf{FINAL EXAMINATION} \\ \small Semester II -- Academic Year 2025--2026}
\fancyfoot[L]{\small (This test has \pageref{LastPage} pages)}
\fancyfoot[R]{\small Page \thepage/\pageref{LastPage}}
\renewcommand{\headrulewidth}{0.4pt}

\setlength{\headheight}{30pt}
\setlength{\parindent}{0pt}
\setlength{\parskip}{6pt}

\begin{document}

\begin{center}
\textbf{\large TEST 2}\\[2pt]
\textbf{Course:} Applied Statistics for Engineering and Scientists 1 \quad
\textbf{Course ID:} STAT451\\[2pt]
\textbf{Duration:} 90 minutes \hspace{1.5cm} \textbf{Materials:} One A4 formula sheet allowed
\end{center}

\vspace{-6pt}
\noindent\textbf{Student name:} \dotfill \quad \textbf{Student ID:} \dotfill \quad \textbf{No.:} \dotfill

\vspace{4pt}
\noindent\textbf{Note:}
\begin{itemize}[leftmargin=*,topsep=2pt,itemsep=0pt]
  \item All calculations must be rounded to \textbf{four} decimal places.
  \item Standard normal $\mathcal{N}(0,1)$ and Student $t$ distribution tables are provided.
\end{itemize}

%--------------------------------------------------------------------
\textbf{Question 1.} (2.0 \emph{points})
The diameter of stainless steel rods manufactured on an extruder is
normally distributed with a mean of $\mu = 0.503$ inch and a standard
deviation of $\sigma = 0.004$ inch. Rods are classified as acceptable
only if their diameter falls in the interval $[0.495, 0.510]$ inch.
\begin{enumerate}[label=(\alph*),leftmargin=*]
  \item Compute the probability that a randomly selected rod is
        \emph{not} acceptable (i.e., is defective).
  \item Among a random sample of $n = 120$ rods, let $X$ denote the
        number of defective rods. Using the normal approximation to
        the binomial distribution (with continuity correction), find
        $P(X \geq 10)$.
  \item A random sample of $n = 25$ rods is drawn; let $\bar{X}$ be
        the sample mean diameter. Determine the probability that
        $\bar{X}$ differs from the process mean by less than
        0.001 inch.
\end{enumerate}

%--------------------------------------------------------------------
\textbf{Question 2.} (2.0 \emph{points})
A random sample of $n = 15$ specimens of a particular type of aluminum
alloy yielded the following yield strengths (in ksi):
\begin{center}
33.8, 35.6, 34.2, 36.1, 33.4, 35.0, 34.7, 36.3,\\
33.9, 35.2, 34.5, 35.9, 33.6, 34.8, 36.0
\end{center}
Assume the yield strengths are approximately normally distributed.
\begin{enumerate}[label=(\alph*),leftmargin=*]
  \item Compute the sample mean $\bar{x}$ and sample standard
        deviation $s$. Then construct a 98\% confidence interval for
        the true mean yield strength $\mu$.
  \item In a quality audit, 186 out of 250 randomly inspected
        aluminum castings passed all surface-defect tests. Construct
        a 95\% confidence interval for the true proportion $p$ of
        castings that pass.
  \item Using the preliminary estimate from part (b), how many
        castings should be inspected to estimate $p$ to within
        0.03 with 95\% confidence?
\end{enumerate}

%--------------------------------------------------------------------
\textbf{Question 3.} (2.0 \emph{points})
The specification sheet for a certain chemical process asserts that
its mean yield is exactly 500 grams per batch. A process engineer
suspects that recent production has drifted from this target. A random
sample of $n = 16$ batches produced a sample mean yield of
$\bar{x} = 493.4$ grams with a sample standard deviation of
$s = 11.2$ grams. Yields are approximately normally distributed.
\begin{enumerate}[label=(\alph*),leftmargin=*]
  \item At the $\alpha = 0.05$ significance level, test the hypothesis
        $H_0: \mu = 500$ against $H_1: \mu \neq 500$. State the
        rejection region, the test statistic, and your conclusion.
  \item In a separate study, the engineer records that 58 out of 400
        randomly selected batches failed to meet the yield target.
        The company had previously claimed the true failure rate is
        at most 12\%. Test this claim at the $\alpha = 0.01$
        significance level using the rejection-region approach.
  \item Compute the $p$-value for part (b) and interpret its meaning.
\end{enumerate}

%--------------------------------------------------------------------
\textbf{Question 4.} (1.5 \emph{points})
Two machines are used to fill plastic bottles with a net volume of
16.0 ounces. The fill volumes are assumed to be normal. To compare the
two machines, quality engineers collect random samples with the
following results:
\begin{center}
\begin{tabular}{lcc}
\toprule
 & Machine 1 & Machine 2 \\
\midrule
Sample size                 & $n_1 = 10$   & $n_2 = 10$   \\
Sample mean                 & $\bar{x}_1 = 16.015$ & $\bar{x}_2 = 16.005$ \\
Sample standard deviation   & $s_1 = 0.030$ & $s_2 = 0.025$ \\
\bottomrule
\end{tabular}
\end{center}
\begin{enumerate}[label=(\alph*),leftmargin=*]
  \item Assuming the two population variances are \emph{equal},
        test at $\alpha = 0.05$ whether there is any difference in the
        mean fill volumes.
  \item Construct a 95\% confidence interval for the difference in
        mean fill volumes, $\mu_1 - \mu_2$.
\end{enumerate}

%--------------------------------------------------------------------
\textbf{Question 5.} (2.5 \emph{points})
A researcher investigates the relationship between the hardness
($y$, Rockwell B) of a copper alloy and the percentage of tin ($x$,
\%) added to the alloy. Data for $n = 8$ specimens were recorded:
\begin{center}
\begin{tabular}{l|cccccccc}
$x$ & 0.5 & 1.0 & 1.5 & 2.0 & 2.5 & 3.0 & 3.5 & 4.0 \\ \hline
$y$ & 41.2 & 45.0 & 48.7 & 52.4 & 55.8 & 59.1 & 62.3 & 65.4 \\
\end{tabular}
\end{center}
\begin{enumerate}[label=(\alph*),leftmargin=*]
  \item Compute the sample correlation coefficient $r$ between
        tin percentage and hardness, and interpret it.
  \item Determine the least-squares regression line
        $\hat{y} = \hat{\beta}_0 + \hat{\beta}_1 x$, and interpret
        both estimated coefficients.
  \item Determine the coefficient of determination $R^{2}$ and
        interpret it in the context of this study.
  \item Predict the hardness when the tin percentage is
        $x = 2.75\%$ and when $x = 6.0\%$. For each prediction,
        state whether the estimate is reliable and justify your
        answer.
\end{enumerate}

\vspace{10pt}
\begin{center}
\textbf{--- END OF TEST 2 ---}
\end{center}

\end{document}
